    \documentclass{article}
    \usepackage{amsmath}
	\usepackage{amssymb}
    \usepackage[minted]{tcolorbox}
    \usepackage{xcolor} % Required for bgcolor in minted
    \definecolor{lightgray}{rgb}{0.9,0.9,0.9} % Define lightgray if not already defined
    \usepackage{fontspec} %fontspec is required for code font
    \setmonofont{JuliaMono}[Extension=.ttf, UprightFont=*-Regular, BoldFont=*-Bold, ItalicFont=*-RegularItalic, BoldItalicFont=*-BoldItalic, Contextuals=Alternate, Scale = 0.8]
    \usepackage{tabularx} % For tables that fit page width
    \usepackage{array} % For better column formatting
    \usepackage{booktabs} % For professional looking tables
	\usepackage{fullpage,graphicx,psfrag,amsfonts,verbatim, url}
    \usepackage{parskip}
    
    \title{\textbf{\textsf{\texttt{PEPit.jl} tutorial: Accelerated inexact forward-backward}}}
	
	\author{John Doe}
    
    \begin{document}
    \maketitle
    
    

\section{Overview}

This tutorial shows how to model and solve a \textbf{performance estimation problem (PEP)} for an \textbf{accelerated inexact forward-backward} method (AIFB), also known as an \emph{inexact accelerated proximal gradient method}.

We consider the composite convex minimization problem

\begin{equation}\label{aifb-problem}
F_\star \triangleq \min_x \left\{F(x) \equiv f(x) + g(x)\right\},
\end{equation}

where:

\begin{itemize}
\item $f$ is convex and $L$-smooth,
\item $g$ is closed, proper, and convex,
\item $\nabla f(\cdot)$ is available exactly,
\item $\mathrm{prox}_{\gamma g}$ is available only \textbf{inexactly}, with a relative accuracy parameter $\zeta \in (0,1)$.
\end{itemize}

Our goal is to compute the smallest possible $\tau(n,L,\zeta)$ such that the guarantee

\begin{equation}\label{aifb-goal}
F(x_n) - F_\star \le \tau(n,L,\zeta)\,\|x_0-x_\star\|^2,
\end{equation}

is valid for the iterate $x_n$ produced by AIFB, under the normalization $\|x_0-x_\star\|^2 \le 1$.

\subsection{The AIFB algorithm}

We implement a simplified instance of [1, Algorithm 3.1] (with the same parameter choices as in the reference example), namely

\begin{equation}\label{aifb-iter}
\begin{aligned}
A_{t+1} &= A_t + \frac{\eta+\sqrt{\eta^2+4\eta A_t}}{2}, \\
y_t &= x_t + \frac{A_{t+1}-A_t}{A_{t+1}}(z_t-x_t), \\
(x_{t+1},v_{t+1}) &\approx_{\varepsilon_t} \left(\mathrm{prox}_{\gamma g}\!\left(y_t-\gamma \nabla f(y_t)\right),\,
\mathrm{prox}_{g^*/\gamma}\!\left(\frac{y_t-\gamma \nabla f(y_t)}{\gamma}\right)\right), \\
\varepsilon_t &= \frac{\zeta^2\gamma^2}{2}\|v_{t+1}+\nabla f(y_t)\|^2, \\
z_{t+1} &= z_t-(A_{t+1}-A_t)\left(v_{t+1}+\nabla f(y_t)\right),
\end{aligned}
\end{equation}

with $\gamma=\tfrac{1}{L}$ and $\eta=(1-\zeta^2)\gamma$.

\subsection{Baseline theoretical guarantee}

We will compare the PEPit-computed bound with the theoretical upper bound from [1, Corollary 3.5],

\begin{equation}\label{aifb-theory}
F(x_n)-F_\star \le \frac{2L}{(1-\zeta^2)n^2}\,\|x_0-x_\star\|^2.
\end{equation}
\section{Setup}

We start by loading \texttt{PEPit.jl} and \texttt{OrderedCollections.jl}.

\begin{itemize}
\item \texttt{PEPit.jl} provides the symbolic objects (\texttt{Point}, \texttt{Expression}, constraints) and the SDP builder/solver (\texttt{solve!}).
\item \texttt{OrderedCollections.jl} provides \texttt{OrderedDict}, which this codebase uses for deterministic parameter passing.
\end{itemize}

\begin{tcblisting}{listing only, minted language=julia, minted options={mathescape=true, breaklines}, colback=lightgray, colframe=black, boxrule=0pt, toprule=2pt, bottomrule=2pt, arc=0pt, outer arc=0pt, left=5pt, right=5pt, top=5pt, bottom=5pt}
using PEPit # Load PEPit.jl (PEP modeling primitives and solver entrypoint).
using OrderedCollections # Use `OrderedDict` for deterministic parameter dictionaries.
\end{tcblisting}


\section{Step 1: Instantiate an empty PEP}

We create an empty \texttt{PEP()} object that we will populate with:

\begin{itemize}
\item declared function classes ($f$ and $g$),
\item initial conditions (here $\|x_0-x_\star\|^2 \le 1$),
\item algorithmic relations (the AIFB iteration),
\item and the performance metric $F(x_n)-F_\star$.
\end{itemize}

\begin{tcblisting}{listing only, minted language=julia, minted options={mathescape=true, breaklines}, colback=lightgray, colframe=black, boxrule=0pt, toprule=2pt, bottomrule=2pt, arc=0pt, outer arc=0pt, left=5pt, right=5pt, top=5pt, bottom=5pt}
problem = PEP() # Create an empty PEP container (lists of functions/points/constraints/metrics).
\end{tcblisting}


\section{Step 2: Choose numerical parameters}

We set:

\begin{itemize}
\item the smoothness constant $L$ of $f$,
\item the relative inexactness parameter $\zeta$,
\item the number of iterations $n$.
\end{itemize}

The reference example uses $L=1.3$, $\zeta=0.45$, and $n=11$.

\begin{tcblisting}{listing only, minted language=julia, minted options={mathescape=true, breaklines}, colback=lightgray, colframe=black, boxrule=0pt, toprule=2pt, bottomrule=2pt, arc=0pt, outer arc=0pt, left=5pt, right=5pt, top=5pt, bottom=5pt}
L = 1.3 # Smoothness constant $L$ for the class of $L$-smooth convex functions.
ζ = 0.45 # Relative inexactness parameter $\zeta \in (0,1)$ controlling the proximal accuracy.
n = 11 # Number of AIFB iterations $n$.
verbose = true # Toggle verbose PEPit/solver output.
\end{tcblisting}


\section{Step 3: Declare the function classes $f$ and $g$}

We declare:

\begin{itemize}
\item $f$ as a smooth convex function with parameter $L$,
\item $g$ as a (possibly nonsmooth) convex closed proper function.
\end{itemize}

Then we form the composite objective $F = f + g$.

\begin{tcblisting}{listing only, minted language=julia, minted options={mathescape=true, breaklines}, colback=lightgray, colframe=black, boxrule=0pt, toprule=2pt, bottomrule=2pt, arc=0pt, outer arc=0pt, left=5pt, right=5pt, top=5pt, bottom=5pt}
f = declare_function!(problem, SmoothConvexFunction, OrderedDict("L" => L); reuse_gradient=true) # Declare $f$: convex and $L$-smooth (unique gradient at each point).
g = declare_function!(problem, ConvexFunction, OrderedDict(); reuse_gradient=false) # Declare $g$: closed, proper, convex (subgradients need not be unique).
F = f + g # Define the composite objective $F(x)=f(x)+g(x)$.
\end{tcblisting}


\section{Step 4:  Model $x_\star$ and impose the initial condition}

In a PEP, we do not search over points directly. Instead, we introduce \textbf{symbolic samples} of the oracles that must be consistent with the function class.

To represent an optimizer $x_\star$, we record a \emph{stationary} sample of the composite objective:

\begin{equation*}
0 \in \partial F(x_\star).
\end{equation*}

Then we introduce the initial point $x_0$ and impose the normalization $\|x_0-x_\star\|^2 \le 1$.

\begin{tcblisting}{listing only, minted language=julia, minted options={mathescape=true, breaklines}, colback=lightgray, colframe=black, boxrule=0pt, toprule=2pt, bottomrule=2pt, arc=0pt, outer arc=0pt, left=5pt, right=5pt, top=5pt, bottom=5pt}
xs = stationary_point!(F) # Create a symbolic minimizer $x_\star$ by enforcing $0 \in \partial F(x_\star)$.
F_star = value!(F, xs) # Define the optimal value $F_\star = F(x_\star)$.
x0 = set_initial_point!(problem) # Create the initial point $x_0$ (a new leaf `Point`).
set_initial_condition!(problem, (x0 - xs)^2 <= 1) # Impose $\|x_0-x_\star\|^2 \le 1$ as the initial condition.
\end{tcblisting}


\section{Step 5: Encode the AIFB iteration (Equation @eq-aifb-iter)}

We now implement the iteration in Equation \eqref{aifb-iter}.

\subsection{About the inexact proximal step in \texttt{PEPit.jl}}

We call \texttt{inexact\_proximal\_step!} with \texttt{opt="PD\_gapI"} to model an \textbf{inexact} proximal computation of $g$ at an input $x_0$.
The primitive returns an \texttt{Expression} \texttt{eps\_var} that represents an upper bound on a primal-dual gap quantity (see \texttt{src/primitive\_steps/inexact\_proximal\_step.jl}).

To match Equation \eqref{aifb-iter}, we add the extra (relative) accuracy constraint

\begin{equation}\label{aifb-accuracy}
\varepsilon_t \le \frac{\zeta^2\gamma^2}{2}\,\|v_{t+1}+\nabla f(y_t)\|^2.
\end{equation}

This is the key modeling step that couples the inexactness level to the algorithm state.

\begin{tcblisting}{listing only, minted language=julia, minted options={mathescape=true, breaklines}, colback=lightgray, colframe=black, boxrule=0pt, toprule=2pt, bottomrule=2pt, arc=0pt, outer arc=0pt, left=5pt, right=5pt, top=5pt, bottom=5pt}
γ = 1 / L # Set $\gamma = 1/L$ (forward step-size).
η = (1 - ζ^2) * γ # Set $\eta = (1-\zeta^2)\gamma$ (parameter used in $A_{t+1}$).
A = 0.0 # Initialize the acceleration scalar $A_0 = 0$.
x = x0 # Initialize the iterate $x_0$.
z = x0 # Initialize the auxiliary iterate $z_0$.
hx = nothing # Placeholder for the final value $g(x_n)$ (filled inside the loop).

for t in 0:(n - 1) # Run $n$ iterations, indexed as $t=0,\dots,n-1$.
    A_next = A + (η + sqrt(η^2 + 4 * η * A)) / 2 # Update $A_{t+1}=A_t+\frac{\eta+\sqrt{\eta^2+4\eta A_t}}{2}$.
    y = x + (1 - A / A_next) * (z - x) # Extrapolation: $y_t=x_t+\frac{A_{t+1}-A_t}{A_{t+1}}(z_t-x_t)$.
    gy = gradient!(f, y) # Gradient query: $g_y=\nabla f(y_t)$.
    x_next, _, hx_next, _, v_next, _, eps_var = inexact_proximal_step!(y - γ * gy, g, γ; opt="PD_gapI") # Inexact prox on $g$ at $y_t-\gamma\nabla f(y_t)$ (returns $x_{t+1}$, $v_{t+1}$, $\varepsilon_t$).
    add_constraint!(g, eps_var <= (ζ * γ)^2 / 2 * (v_next + gy)^2) # Enforce Equation @eq-aifb-accuracy: $\varepsilon_t \le \frac{\zeta^2\gamma^2}{2}\|v_{t+1}+\nabla f(y_t)\|^2$.
    z = z - (A_next - A) * (v_next + gy) # Update $z_{t+1}=z_t-(A_{t+1}-A_t)(v_{t+1}+\nabla f(y_t))$.
    x = x_next # Commit the new iterate $x_{t+1}$.
    hx = hx_next # Store $g(x_{t+1})$ to later build $F(x_n)=f(x_n)+g(x_n)$.
    A = A_next # Commit the scalar update $A_t \leftarrow A_{t+1}$.
end # End the AIFB loop.
\end{tcblisting}


\section{Step 6: Define the performance metric and solve the PEP}

The performance metric is the final objective error:

\begin{equation*}
F(x_n) - F_\star = f(x_n) + g(x_n) - F_\star.
\end{equation*}

In PEPit, we encode this quantity as an \texttt{Expression} and pass it to \texttt{set\_performance\_metric!}.
Then \texttt{solve!} constructs and solves the SDP and returns the worst-case value.

\begin{tcblisting}{listing only, minted language=julia, minted options={mathescape=true, breaklines}, colback=lightgray, colframe=black, boxrule=0pt, toprule=2pt, bottomrule=2pt, arc=0pt, outer arc=0pt, left=5pt, right=5pt, top=5pt, bottom=5pt}
set_performance_metric!(problem, value!(f, x) + hx - F_star) # Set the performance metric to $F(x_n)-F_\star$.
pepit_tau = solve!(problem; verbose=verbose) # Solve the SDP and return the worst-case value $\tau(n,L,\zeta)$.
\end{tcblisting}


\section{Step 7: Compare with the theoretical upper bound (Equation @eq-aifb-theory)}

The theory predicts the upper bound in Equation \eqref{aifb-theory}:

\begin{equation*}
\tau_{\mathrm{th}}(n,L,\zeta) = \frac{2L}{(1-\zeta^2)n^2}.
\end{equation*}

We compute it and print both values.

\begin{tcblisting}{listing only, minted language=julia, minted options={mathescape=true, breaklines}, colback=lightgray, colframe=black, boxrule=0pt, toprule=2pt, bottomrule=2pt, arc=0pt, outer arc=0pt, left=5pt, right=5pt, top=5pt, bottom=5pt}
theoretical_tau = 2 * L / (1 - ζ^2) / n^2 # Compute $\tau_{\mathrm{th}}=\frac{2L}{(1-\zeta^2)n^2}$ from Equation @eq-aifb-theory.

@info "📼 Accelerated inexact forward-backward (AIFB) ***" # Print a short header.

@info "💻 PEPit guarantee:\t F(x_n)-F_* <= $(round(pepit_tau, digits=7)) ||x_0 - x_*||^2" # Report the PEPit-computed value.

@info "📰 Theoretical bound:\t F(x_n)-F_* <= $(round(theoretical_tau, digits=7)) ||x_0 - x_*||^2" # Report the theoretical comparison bound.
\end{tcblisting}


\subsection{Expected output (for the reference parameters)}

Assuming everything went well, with $L=1.3$, $\zeta=0.45$, and $n=11$, the reference run reports approximately:

\begin{tcblisting}{listing only, minted language=julia, minted options={mathescape=true, breaklines}, colback=lightgray, colframe=black, boxrule=0pt, toprule=2pt, bottomrule=2pt, arc=0pt, outer arc=0pt, left=5pt, right=5pt, top=5pt, bottom=5pt}
PEPit guarantee:      F(x_n)-F_* <= 0.0187341 ||x_0 - x_*||^2
Theoretical bound:    F(x_n)-F_* <= 0.0269437 ||x_0 - x_*||^2
\end{tcblisting}


Our numerical values may differ slightly depending on the SDP solver and tolerances. 
\section{Reference}

[1] M. Barre, A. Taylor, F. Bach (2021).
\emph{A note on approximate accelerated forward-backward methods with absolute and relative errors, and possibly strongly convex objectives.}
arXiv:2106.15536v2.
URL: https://arxiv.org/pdf/2106.15536v2.pdf


\end{document}
